\begin{textsample}{POS Dim 4 – 2018\_c6 – Score 12.00 – 2018\_c6/enade\_2018\_des\_d0d272ca78.txt}  \label{ex:f4_pos_019}
CAPÍTULO 6 ANÁLISE TÉCNICA DA PROVA Este capítulo tem por objetivo apresentar o desempenho dos estudantes concluintes de Design no Enade/2018.
Para isso, foram calculadas as estatísticas básicas da prova em seu todo ( seção 6.1.1 ), bem como as estatísticas relacionadas aos Componentes de Formação Geral ( seção 6.1.2 ) e de Conhecimento Específico da Área ( seção 6.1.3 ).
Dadas as suas características, foram analisadas, em separado, as questões objetivas ( seção 6.2 ) e as questões discursivas ( seção 6.3 ).
Para as questões objetivas, foram disponibilizados os índices de facilidade e de discriminação Ponto-Bisserial, também em separado, para os Componentes de Formação Geral ( seção 6.2.1 ) e de Conhecimento Específico ( seção 6.2.2 ).
De cada componente, uma das questões foi escolhida para exemplificar a análise gráfica, relacionando as alternativas escolhidas pelos estudantes ( inclusive o gabarito ) com o número de acertos no componente.
No Anexo I, apresenta -se a íntegra da análise gráfica para todas as questões objetivas.
Para cada uma das questões discursivas, os conteúdos dos tipos mais comuns de respostas dos estudantes são apresentados e comparados com o padrão de respostas esperado ( ver Anexo VIII com o padrão de respostas ).
Tomando -se como base as duas questões discursivas do Componente de Formação Geral, na seção 6.3.1.6 são apresentados comentários sobre a correção das respostas em relação à Língua Portuguesa.
Nas tabelas, constam as seguintes estatísticas das notas²⁰ : média do desempenho na prova, erro padrão da média, desvio padrão, nota mínima, mediana e nota máxima para cada um de seus componentes.
Tais estatísticas contemplam o total de estudantes concluintes da Área de Design, inscritos e presentes à prova do Enade/2018, tendo em vista agregações, ou por Grandes Regiões e o país como um todo, ou por Categoria Administrativa, Organização Acadêmica da IES e Modalidade de Ensino.
Em relação aos gráficos de distribuição de notas, o intervalo considerado foi de 10 unidades, aberto à esquerda e fechado à direita, com exceção do primeiro intervalo, [ 0 ; 10 ], fechado em ambos os extremos.
Para os gráficos de distribuição das notas das questões discursivas, foram consideradas mais duas categorias : questão em branco²¹ e nota zero. ²⁰ Uma definição dessas estatísticas pode ser encontrada no Glossário. 6.1 ESTATÍSTICAS BÁSICAS DA PROVA Esta seção apresenta estatísticas selecionadas e histogramas da nota geral ( 6.1.1 ) e de cada componente : Formação Geral ( 6.1.2 ) e Conhecimento Específico ( 6.1.3 ).
São também apresentadas estatísticas selecionadas de subpopulações, caracterizadas por Grande Região, Categoria Administrativa, Organização Acadêmica e Modalidade de Ensino. 6.1.1 Estatísticas Básicas Gerais Na Tabela 6.1, são apresentadas as Estatísticas Básicas da prova, por Grande Região, dos estudantes concluintes de Design.
A Média das notas da prova, como um todo ( nas seções seguintes serão analisados os Componentes de Formação Geral e de Conhecimento Específico ), foi 50,8, sendo que os alunos da região Centro-Oeste obtiveram a Média mais alta ( 54,9 ), e os da região Norte obtiveram a Média mais baixa ( 47,0 ).
As demais médias foram : 50,1, na região Nordeste, 50,5, na região Sudeste e 51,4, na região Sul.
O Desvio padrão para o Brasil, como um todo, foi 13,0, sendo o maior Desvio padrão encontrado na região Centro-Oeste ( 14,1 ), e o menor, nas regiões Norte, Nordeste e Sul ( 12,9 ), indicando uma dispersão um pouco menor das notas dessas três últimas regiões.
A região que obteve a maior nota Máxima foi a Sul ( 89,2 ), ao passo que a região que atingiu a menor nota Máxima foi a Norte ( 74,9 ).
A Mediana do Brasil, como um todo, foi 51,9, sendo a maior Mediana obtida na região Centro-Oeste ( 57,5 ), e a menor, obtida na Norte ( 49,4 ).
A nota Mínima do Brasil foi obtida nas regiões Nordeste, Sudeste e Sul, onde essa foi zero.
A Mínima da região Norte foi 4,4, e a da região foi 14,0, a maior nota Mínima dentre as cinco Grandes Regiões.
Considerando -se as notas segundo Grande Região, observa -se que existe diferença estatisticamente significativa ao nível de 95 \% entre a maior Média, obtida na região Centro-Oeste, e as médias das demais regiões²². ²² Todas as comparações deste capítulo utilizam os intervalos de 95 \%.
Os erros-padrão da média que possibilitam os testes estão disponíveis nas tabelas.
O comportamento das notas dos estudantes de todo o Brasil pode ser observado no Gráfico 6.1, que apresenta um histograma com a distribuição das mesmas : uma distribuição unimodal com a moda no intervalo ( 50 ; 60 ], seguido do intervalo ( 40 ; 50 ].
Na Tabela 6.2, são apresentadas informações referentes à Média da nota final, desagregadas por Categoria Administrativa, por Organização Acadêmica e por Modalidade de Ensino.
A Média das notas da prova, como um todo, foi 50,8.
Em relação à Categoria Administrativa, os estudantes das IES Públicas obtiveram Média mais alta ( 55,1 ) que a Média nacional, e os das IES Privadas obtiveram Média mais baixa ( 47,9 ).
Observa -se que existe diferença estatisticamente significativa entre as médias das notas das IES Públicas e as das Privadas.
A diferença entre as médias das regiões Centro-Oeste e Norte ( 7,9 ), a maior e a menor Média, é superior à diferença entre a das IES Públicas e a das Privadas ( 7,2 ), caracterizando -se uma maior diversidade administrativa do que regional.
O Desvio padrão para as IES Públicas ( 12,5 ) e para IES Privadas ( 12,6 ) foi inferior ao do Brasil, como um todo ( 13,0 ).
No tocante à Organização Acadêmica, as Universidades obtiveram Média mais alta que a nacional ( 52,0 e 50,8, respectivamente ).
A Média dos Centros Universitários, Faculdades e CEFET/IFET foram menores do que a nacional ( 47,3, 47,2 e 48,7, respectivamente ).
Constata -se que existe diferença, estatisticamente significativa, ao nível de 95 \% entre a Média mais alta, 52,0 das Universidades, e as duas mais baixas, obtidas em Centros Universitários e Faculdades.
Como não existem cursos de Design na Modalidade Educação a Distância, os dados referentes à Educação Presencial são idênticos ao Total do Brasil como um todo da tabela anterior. 6.1.2 Estatísticas Básicas no Componente de Formação Geral Na Tabela 6.3, são apresentadas as Estatísticas Básicas em relação ao componente da prova que avalia a Formação Geral dos estudantes concluintes.
Os alunos de todo o Brasil obtiveram desempenho médio de 52,9.
Quanto à variabilidade, o Desvio padrão das notas dos estudantes do Brasil, como um todo, foi 18,4.
A maior Média foi obtida na região Centro-Oeste ( 57,6 ), e a menor, na região Norte ( 47,8 ).
As demais médias foram : 51,5, na região Nordeste, 53,6, na região Sudeste e 52,5, na região Sul.
Já o maior Desvio padrão foi obtido nas regiões Norte e Centro-Oeste ( 18,8 ), e o menor, na região Nordeste ( 18,1 ).
Os demais desvios padrões foram iguais a 18,4, tanto para a região Sudeste como para a região Sul.
A maior nota no Componente de Formação Geral da prova do Enade/2018 foi 96,4, obtida por, pelo menos, um aluno nas regiões Sudeste e Sul.
A menor nota Máxima foi obtida na região Norte ( 88,3 ).
A Mediana do Brasil, como um todo, foi 53,3, sendo a menor Mediana encontrada na região Norte ( 47,6 ), e a maior, na região Centro-Oeste ( 60,3 ).
A nota Mínima nessa parte foi zero em quase todas as regiões, com exceção das regiões Norte, cuja nota Mínima foi 5,2, e Centro-Oeste, cuja nota Mínima foi 7,5.
Considerando -se as notas, segundo Grande Região, observa -se que existe diferença estatisticamente significativa entre a Média das notas do Componente de Formação Geral obtida na região Centro-Oeste, a maior Média, e a Média das demais regiões. nas regiões Norte e Centro-Oeste ( 18,8 ), e o menor, na região Nordeste ( 18,1 ).
Os demais desvios padrões foram iguais a 18,4, tanto para a região Sudeste como para a região Sul.
A maior nota no Componente de Formação Geral da prova do Enade/2018 foi 96,4, obtida por, pelo menos, um aluno nas regiões Sudeste e Sul.
A menor nota Máxima foi obtida na região Norte ( 88,3 ).
A Mediana do Brasil, como um todo, foi 53,3, sendo a menor Mediana encontrada na região Norte ( 47,6 ), e a maior, na região Centro-Oeste ( 60,3 ).
A nota Mínima nessa parte foi zero em quase todas as regiões, com exceção das regiões Norte, cuja nota Mínima foi 5,2, e Centro-Oeste, cuja nota Mínima foi 7,5.
Considerando -se as notas, segundo Grande Região, observa -se que existe diferença estatisticamente significativa entre a Média das notas do Componente de Formação Geral obtida na região Centro-Oeste, a maior Média, e a Média das demais regiões.
No Gráfico 6.2, é possível se verificar a avaliação do desempenho dos estudantes do Componente de Formação Geral, a partir do histograma da distribuição das notas correspondentes.
A distribuição é unimodal, com moda em ( 50 ; 60 ], o mesmo intervalo da distribuição de notas da prova, como um todo ( Gráfico 6.1 ).
Nota -se, ainda, que, no Gráfico 6.2, as notas apresentam maior dispersão do que as no Gráfico 6.1 ( distribuição das notas da prova ), confirmado pela comparação dos desvios padrões : 13,0 para a nota da prova, como um todo, e 18,4 para o Componente de Formação Geral.
Na Tabela 6.4, são apresentadas as informações referentes ao desempenho dos concluintes do Componente de Formação Geral, em diferentes agregações : Categoria Administrativa e Organização Acadêmica.
Observa -se que existe diferença estatisticamente significativa entre as médias dos tipos de Categoria Administrativa : a maior Média obtida por estudantes de IES Públicas ( 58,5 ), e a menor, pelos de IES Privadas ( 49,3 ).
Nota -se que também há diferença estatisticamente significativa entre a maior das médias dos quatro tipos de Organização Acadêmica e as duas maiores a seguir, o valor maior para as Universidades ( 54,7 ) e os dois maiores a seguir, de Faculdades ( 48,6 ) e de Centros Universitários ( 47,2 ). 6.1.3 - Estatísticas Básicas do Componente de Conhecimento Específico Na Tabela 6.5, são apresentadas as Estatísticas Básicas referentes ao Componente de Conhecimento Específico da área de Design.
A Média do desempenho dos alunos do Brasil, como um todo, foi 50,0.
A maior Média foi obtida na região Centro-Oeste ( 54,0 ), e a menor, na região Norte ( 46,7 ).
As demais médias foram : 49,6, na região Norte, 49,4, na região Sudeste e 51,1, na região Sul.
Quanto à variabilidade das notas, o Desvio padrão do Brasil, como um todo, foi 13,2, sendo o maior Desvio padrão observado na região Centro-Oeste ( 14,3 ), e o menor, na região Sul ( 12,9 ).
Os demais desvios foram : 13,6, na região Norte e 13,2, nas regiões Nordeste e Sudeste.
A Mediana das notas dos estudantes de todo o Brasil foi 51,0.
A maior Mediana ocorreu na região Centro-Oeste ( 55,9 ), e a menor, na região Norte ( 47,6 ).
As demais medianas foram : 50,7 na região Nordeste, 50,9, na região Sudeste e 52,2, na região Sul.
A nota Máxima do Brasil, como um todo, foi 89,9, sendo obtida por, pelo menos, um aluno, na região Sul.
As demais notas máximas foram : 72,9, na região Norte, a menor nota Máxima dentre elas, 78,8, na região Nordeste, 86,5, na região Sudeste e 84,4, na região Centro-Oeste.
A nota Mínima foi zero em quase todas as regiões, com exceção da região Norte, que teve nota Mínima 3,4, e da região Centro-Oeste, com nota Mínima de 13,6.
Observa -se que existe diferença estatisticamente significativa entre a maior Média das notas do Componente de Conhecimento Específico, obtida na região Centro-Oeste ( 54,0 ) e as demais médias das regiões.
Assim como nos Gráficos 6.1 e 6.2, no Gráfico 6.3, apresentado a seguir, é possível ser feita uma avaliação do desempenho de concluintes de Design em relação ao Componente de Conhecimento Específico, com um histograma da distribuição das notas correspondentes.
Esta também é uma distribuição unimodal, e o grupo modal é o ( 50 ; 60 ], o mesmo grupo modal da prova, como um todo, e da Formação Geral.
Na Tabela 6.6, é apresentada uma comparação dos resultados em relação à Categoria Administrativa e à Organização Acadêmica, agora levando -se em conta o desempenho de alunos do Componente de Conhecimento Específico da prova.
No que se refere à Organização Acadêmica, a maior Média foi a das Universidades ( 51,0 ), vindo a seguir a dos CEFET/IFET ( 49,4 ), a dos Centros Universitários ( 47,3 ) e, por fim, a das Faculdades ( 46,7 ).
O maior Desvio padrão, igual ao valor para o Brasil, como um todo, foi o das Universidades e Faculdades ( 13,2 ).
As Universidades obtiveram a maior nota Máxima ( 89,9 ).
Centros Universitários e Faculdades obtiveram a mesma nota Máxima 78,4 ; e os CEFET/IFET, nota Máxima 78,6.
As medianas foram : 52,4, nas Universidades, a maior delas ; 48,0, nos Centros Universitários ; 47,6, na Faculdades, a menor delas, e 50,7, nos, CEFET/IFET.
A nota Mínima foi zero nas Universidades, 6,8, nos Centros Universitários, 3,4, nas Faculdades e 31,2, nos CEFET/IFET, a maior nota Mínima.
Observa -se que existe diferença estatisticamente significativa ao nível de 95 \% no Componente de Conhecimento Específico entre as notas das Universidades, a maior Média, e as dos Centros Universitários e das Faculdades, as duas menores.
Quanto à Categoria Administrativa, observa -se um comportamento semelhante àquele da parte de Formação Geral e à prova, como um todo, ou seja, existe diferença estatisticamente significativa entre as médias das IES Públicas ( 53,9 ) e as das IES Privadas ( 47,5 ).
Nesse caso também, a maior Média foi obtida por alunos de IES Públicas de ensino. 6.2 ANÁLISE DAS QUESTÕES OBJETIVAS Esta seção apresenta estatísticas selecionadas e histogramas das Questões Objetivas de Formação Geral ( 6.2.1 ) e de Conhecimento Específico ( 6.2.2 ).
São também apresentadas e comparadas as médias das subpopulações caracterizadas por Grande Região. 6.2.1 Componente de Formação Geral Na Tabela 6.7, são apresentadas as Estatísticas Básicas relativas às oito questões objetivas do componente da prova que abrange a Formação Geral dos estudantes.
A Média do Brasil foi 60,6.
A menor Média foi encontrada na região Norte ( 55,4 ), e a maior, na região Centro-Oeste ( 67,6 ).
As demais médias foram : 59,0, na região Nordeste, 61,0, na região Sudeste e 60,6, na região Sul.
O Desvio padrão do Brasil foi 23,6, sendo o maior Desvio padrão encontrado na região Centro-Oeste ( 24,1 ), e o menor, na região Norte ( 22,7 ).
Os demais desvios foram : 23,6, na região Nordeste, 23,4, na região Sudeste e 23,9, na região Sul.
As medianas do Brasil, como um todo, e de três das cinco regiões foram 62,5.
A Mediana da região Norte foi 50,0, menor, e a da região Centro-Oeste foi 75,0, maior que a nacional.
A nota Máxima 100,0 e a nota Mínima zero foram iguais para todas as Grandes Regiões.
Na Tabela 6.8, são apresentados o Índice de Facilidade e o Índice de Discriminação ( Ponto-Bisserial ) para cada uma das questões objetivas do Componente de Formação Geral.
Quanto ao Índice de Facilidade, foram usadas as seguintes cores para diferenciar o nível de dificuldade da questão :  Azul para as questões classificadas com índice Muito fácil ( \&gt ; =0,86 ), verde para as questões classificadas com índice Fácil ( 0,61 a 0,85 ), amarelo para as questões classificadas com Médio ( 0,41 a 0,60 ), vermelho para as questões classificadas com Difícil ( 0,16 a 0,40 ) e roxo para as questões classificadas com Muito difícil ( \&lt ; =0,15 ).
Já quanto ao Índice de Discriminação, foram usadas as seguintes cores para qualificar a questão :  As questões classificadas com índice Fraco receberam a cor vermelha ( \&lt ; =0,19 ), as classificadas com Médio receberam a cor amarela ( 0,20 a 0,29 ), as classificadas com Bom receberam a cor verde ( 0,30 a 0,39 ) e as classificadas com Muito bom ( \&gt ; =0,40 ) receberam a cor azul.
As questões objetivas do Componente de Formação Geral, segundo o Índice de Facilidade, foram assim avaliadas : das oito questões, nenhuma teve o Índice de Facilidade As medianas do Brasil, como um todo, e de três das cinco regiões foram 62,5.
A Mediana da região Norte foi 50,0, menor, e a da região Centro-Oeste foi 75,0, maior que a nacional.
A nota Máxima 100,0 e a nota Mínima zero foram iguais para todas as Grandes Regiões.
Na Tabela 6.8, são apresentados o Índice de Facilidade e o Índice de Discriminação ( Ponto-Bisserial ) para cada uma das questões objetivas do Componente de Formação Geral.
Quanto ao Índice de Facilidade, foram usadas as seguintes cores para diferenciar o nível de dificuldade da questão : • Azul para as questões classificadas com índice Muito fácil ( \&gt ; =0,86 ), verde para as questões classificadas com índice Fácil ( 0,61 a 0,85 ), amarelo para as questões classificadas com Médio ( 0,41 a 0,60 ), vermelho para as questões classificadas com Difícil ( 0,16 a 0,40 ) e roxo para as questões classificadas com Muito difícil ( \&lt ; =0,15 ).
Já quanto ao Índice de Discriminação, foram usadas as seguintes cores para qualificar a questão : • As questões classificadas com índice Fraco receberam a cor vermelha ( \&lt ; =0,19 ), as classificadas com Médio receberam a cor amarela ( 0,20 a 0,29 ), as classificadas com Bom receberam a cor verde ( 0,30 a 0,39 ) e as classificadas com Muito bom ( \&gt ; =0,40 ) receberam a cor azul.
As questões objetivas do Componente de Formação Geral, segundo o Índice de Facilidade, foram assim avaliadas : das oito questões, nenhuma teve o Índice de Facilidade classificado como Muito fácil, e quatro questões foram classificadas como Fácil, por terem índice de acertos entre 0,61 e 0,74.
Outras quatro questões foram consideradas com Índice de Dificuldade Médio, situando -se no intervalo entre 0,41 e 0,60 do Índice de Facilidade, ou seja, houve entre 41,0 \% e 60,0 \% de acertos.
Nenhuma questão foi classificada na categoria Difícil, situando -se no intervalo entre 0,16 e 0,40, ou apresentou menos de 15 \% de acertos, razão pela qual também não houve questão classificada como Muito difícil.
O Índice de Facilidade variou de 0,43 a 0,74, e o de Discriminação, de 0,45 a 0,54.
Na Tabela 6.9, é apresentada a distribuição das questões, levando -se em conta, simultaneamente, a classificação dos dois índices.
Todas as oito questões foram classificadas com Índice de Discriminação Muito bom e figuraram entre dois níveis de dificuldade Fácil e Médio : as quatro classificadas na categoria Fácil foram as questões 1,2, 3 e 8 ; as quatro, na categoria Médio foram as questões 4, 5, 6 e 7.
O máximo de acertos foi alcançado pela questão 8 com um Índice de Facilidade de 0,74.
Em particular, as questões 4 e 5 foram as que apresentaram maior poder discriminatório, com índice 0,54, cada, e em termos de facilidade, com proporções de 0,43 e 0,59 acertos, respectivamente.
A questão de número 4 foi a que apresentou o menor Índice de Facilidade 0,43, ou seja, menos da metade dos estudantes conseguiu resolvê -la, dentro do universo de participantes.
No Gráfico 6.4, para exemplificar, analisa -se o comportamento da questão de número 4 de Formação Geral.
Trata -se de uma questão considerada Médio, em relação à facilidade e a que obteve o maior Índice de Discriminação dessa parte da prova, juntamente com a questão 5.
Neste gráfico, cada uma das cinco curvas representa o percentual de respostas em determinada alternativa da questão, em função do número de acertos dos estudantes nessa parte da prova ( Formação Geral/Múltipla Escolha ), antes de possíveis eliminações pelo critério do Ponto-Bisserial.
A curva em verde corresponde à alternativa B, a alternativa correta para essa questão.
Observa -se que dentre os estudantes com menor número de acertos, nessa parte do exame, a situação mais frequente foi a escolha da alternativa C ( em preto ), incorreta.
Por exemplo, dentre os estudantes que acertaram três questões, 33,1 \% escolheram a alternativa C ( em preto ), 18,9 \% escolheram a alternativa E ( em vermelho ), 16,2 \% escolheram a alternativa B ( em verde ), a alternativa correta, 15,8 \% escolheram a alternativa A ( em azul ), e 15,7 \%, a D ( em laranja ).
Entre os que acertaram três respostas entre as questões de múltipla escolha de Formação Geral, 0,1 \% deixou a questão em branco e 0,1 \% marcou mais de uma alternativa, invalidando a questão.
À medida que o número de acertos aumenta, indicando desempenho melhor nessa parte da prova, aumenta, concomitantemente, a proporção de estudantes que selecionaram a alternativa correta B, atingindo 100 \% para os estudantes com oito acertos.
Essa análise permite verificar como a questão discriminou os grupos de desempenho, justificando -se o alto índice obtido na questão.
Cumpre notar que não é possível inferir desse gráfico nem o índice de facilidade ( que seria uma média da proporção ponderada pela quantidade de alunos com cada uma das notas ), nem o índice de discriminação Ponto-Bisserial, por razão equivalente.
No caso extremo, no qual a grande concentração dos acertos dos alunos fosse abaixo de quatro, o índice de facilidade seria obrigatoriamente abaixo de 30 \% ( neste exemplo ).
Caso a concentração fosse em seis acertos ou mais, o índice seria obrigatoriamente acima de 50 \%.
Os gráficos relativos às demais questões de Formação Geral constam do Anexo I. 6.2.2 Componente de Conhecimento Específico Na Tabela 6.10, são apresentadas as Estatísticas Básicas em relação às questões objetivas do Componente de Conhecimento Específico da prova de Design por Grande Região.
Como será discutido ainda nesta seção, três questões objetivas do Componente de Conhecimentos Específicos não foram usadas no cômputo das notas por terem sido descartadas pelo critério do Ponto-Bisserial.
Assim, as notas, foram calculadas com base em 24 das 27 questões objetivas de Conhecimentos Específicos.
A Média do Brasil desse componente foi 53,7.
A menor Média foi observada na região Norte ( 51,0 ), e a maior, na região Centro-Oeste ( 57,1 ).
O Desvio padrão de todo o Brasil foi 14,2, sendo o menor Desvio padrão encontrado na região Sul ( 14,0 ), e o maior, na região Norte ( 14,7 ).
A Mediana de todo o Brasil foi 56,0, o mesmo valor da Mediana encontrada nas regiões Nordeste, Sudeste e Sul.
Na região Norte, a Mediana foi menor : 52,0, e, na região Centro-Oeste, a Mediana foi maior que a nacional : 60,0.
A nota Máxima da prova foi 92,0, obtida nas questões objetivas do Componente de Conhecimento Específico, por, pelo menos, um aluno Na Tabela 6.11, são apresentados os Índices de Facilidade e Discriminação ( Ponto-Bisserial ) das questões objetivas do Componente de Conhecimento Específico, para os estudantes de Design.
Para facilitar a diferenciação das questões, serão usadas as mesmas cores da Tabela 6.8 para as diferentes classificações dos Índices de Facilidade e de Discriminação.
A partir do Índice de Facilidade obtido, pode -se concluir que mais de um terço das questões objetivas da prova foram consideradas, pelo menos, Difícil : das 27 questões válidas, dez foram classificadas como Difícil ou como Muito difícil, oito e duas questões, respectivamente.
Foram classificadas como Fácil um total de oito questões, e outras nove foram consideradas como Médio.
Já quanto aos índices de discriminação das questões objetivas do Componente de Conhecimento Específico da prova, obtém -se como resultado a seguinte classificação : oito das 27 questões foram consideradas boas, enquanto seis delas tiveram Índice de Discriminação Muito bom.
Assim, para um pouco mais da metade das questões – 14 em 27 – os Índices de Discriminação foram Bom ou Muito bom.
Dentre as demais, dez delas foram classificadas como Médio, e outras três, como Fraco, sendo 13, por conseguinte, a quantidade de questões nos dois patamares mais baixos de discriminação.
Constata -se, assim, que a prova – no que se refere ao Componente de Conhecimento Específico – possuía média capacidade de discriminar dentre aqueles que dominam ou não o conteúdo. nas regiões Sudeste e Sul.
A nota Mínima foi zero em três regiões : Nordeste, Sudeste e Sul.
Na região Norte, a nota Mínima foi 4,0 e, na região Centro-Oeste, foi 16,0, a maior dentre as regiões.
Na Tabela 6.11, são apresentados os Índices de Facilidade e Discriminação ( Ponto-Bisserial ) das questões objetivas do Componente de Conhecimento Específico, para os estudantes de Design.
Para facilitar a diferenciação das questões, serão usadas as mesmas cores da Tabela 6.8 para as diferentes classificações dos Índices de Facilidade e de Discriminação.
A partir do Índice de Facilidade obtido, pode -se concluir que mais de um terço das questões objetivas da prova foram consideradas, pelo menos, Difícil : das 27 questões válidas, dez foram classificadas como Difícil ou como Muito difícil, oito e duas questões, respectivamente.
Foram classificadas como Fácil um total de oito questões, e outras nove foram consideradas como Médio.
Já quanto aos índices de discriminação das questões objetivas do Componente de Conhecimento Específico da prova, obtém -se como resultado a seguinte classificação : oito das 27 questões foram consideradas boas, enquanto seis delas tiveram Índice de Discriminação Muito bom.
Assim, para um pouco mais da metade das questões – 14 em 27 – os Índices de Discriminação foram Bom ou Muito bom.
Dentre as demais, dez delas foram classificadas como Médio, e outras três, como Fraco, sendo 13, por conseguinte, a quantidade de questões nos dois patamares mais baixos de discriminação.
Constata -se, assim, que a prova – no que se refere ao Componente de Conhecimento Específico – possuía média capacidade de discriminar dentre aqueles que dominam ou não o conteúdo.
Na Tabela 6.12, é apresentada a distribuição das questões levando -se em conta, simultaneamente, a classificação dos dois índices.
Seis questões obtiveram Índice de Discriminação classificado como Muito bom, sendo uma questão classificada, quanto ao Índice de Facilidade, como Médio e cinco questões, como Fácil.
Das oito questões classificadas como Bom quanto ao Índice de Discriminação, seis foram classificadas, quanto ao Índice de Facilidade, como Médio e duas, como Fácil.
Por sua vez, as dez questões classificadas como Médio, quanto ao Índice de Discriminação, foram classificadas como Difícil ( 8 ), Médio ( 1 ) e Fácil ( 1 ).
E, finalmente, três questões foram classificadas, quanto ao Índice de Discriminação como Fraco, sendo : duas Muito difícil e uma Médio ; quanto ao Índice de Facilidade.
Dentre as questões que alcançaram os maiores índices de discriminação, uma delas foi classificada com Índice de Discriminação Muito bom, a de número 24, seu índice foi 0,48, e, quanto ao índice de Facilidade, foi classificada como Fácil, com 78 \% dos estudantes marcando a opção correta.
A questão de número 26 foi a mais difícil dentre as 27 questões específicas válidas, com baixo Índice de Facilidade, apenas 8 \% de acertos.
Essa questão apresentou poder discriminatório muito baixo, -0,01, o que comprova ter sido a mais difícil para os estudantes.
Destaca -se, também, a questão 19, com Índice de Facilidade 0,14, o que, em termos percentuais, corresponde a 14 \% de estudantes que responderam acertadamente.
Já 0,08 foi o seu Índice de Discriminação.
Tais questões foram, portanto, pelo critério Ponto-Bisserial, consideradas inadequadas.
Por isso, as questões 19 e 26, além da questão 17, foram eliminadas do cômputo da nota final.
A título de exemplo das análises do comportamento das questões objetivas, no Gráfico 6.5, analisa -se a questão 24 do Componente de Conhecimento Específico.
Essa questão foi considerada, pelas respostas dos estudantes avaliados na prova, como uma questão com nível de facilidade Fácil, 0,78, ou seja, 78 \% dos estudantes assinalaram acertadamente a opção A, correspondente ao gabarito.
Como já comentado, seu Índice de Discriminação foi igual a 0,48, classificado como Muito bom, o maior valor de discriminação.
Neste gráfico, cada uma das cinco curvas representa o percentual de respostas em determinada alternativa da questão 24, em função do número de acertos dos estudantes nessa parte da prova, antes de possíveis eliminações de questões pelo critério do Ponto-Bisserial.
A alternativa correta A, representada no gráfico pela curva em azul, foi escolhida em maiores proporções pelos alunos com desempenho melhor nessa parte da prova.
Já as alternativas incorretas, também denominadas distratores, foram selecionadas, principalmente, por aqueles com notas mais baixas.
Observa -se que a soma não é 100 \%, por causa das questões não respondidas ou com mais de uma opção marcada.
Aqueles com nota zero, na sua totalidade, deixaram esta questão em branco ou marcaram mais de uma alternativa, comportamento considerado inválido.
A proporção de alunos que selecionaram a resposta correta A aumenta gradativamente, chegando a atingir 100 \% para 22 acertos ou ou mais, enquanto a proporção dos que escolheram alternativas incorretas decai, a partir de dois ou três acertos, em função do número de acertos nessa parte da prova.
Os gráficos relativos às demais questões do Conhecimento Específico constam do Anexo I. 6.3 ANÁLISE DAS QUESTÕES DISCURSIVAS Esta seção apresenta estatísticas selecionadas e histogramas das Questões Discursivas de Formação Geral ( 6.3.1 ) e Conhecimento Específico ( 6.3.2 ).
São também apresentadas e comparadas as médias de subpopulações, caracterizadas por Grande Região. 6.3.1 Componente de Formação Geral As análises dos resultados de desempenho dos estudantes de Design, nas duas questões discursivas relativas à Formação Geral, encontram -se na Tabela 6.13 e no Gráfico 6.6.
Na Tabela 6.13, observa -se que a nota Média nesse conjunto de questões foi inferior à obtida nas objetivas.
Os estudantes de todo Brasil obtiveram, em Formação Geral, Média 60,6, nas questões objetivas e 41,4, nas questões discursivas.
Pode -se notar também que o Desvio padrão nesse conjunto de questões foi menor do que o obtido nas objetivas : 23,6 nas questões objetivas e 20,6, nas questões discursivas.
A maior Média foi obtida na região Centro-Oeste ( 42,7 ), e a menor, na região Norte ( 36,3 ).
A Mediana de todo o Brasil, neste componente, foi 43,0, a mesma sendo obtida nas regiões Nordeste e Centro-Oeste.
Nas regiões Norte e Sul, a Mediana foi menor, respectivamente 38,5 e 41,5, e, na região Sudeste ( 44,5 ), foi maior que o valor para o Brasil.
A nota Máxima ( 97,5 ) foi obtida na região Sul, sendo 89,5 a Máxima na região Norte ; 90,5, na Nordeste ; 94,5, na Sudeste ; e 88,5, na região Centro-Oeste.
A nota Mínima ( 0,0 ) foi a mesma em todas as regiões do Brasil.
No Gráfico 6.6, está representada a distribuição das notas nas questões discursivas do Componente de Formação Geral.
A moda dessa distribuição ocorre no intervalo ( 40 ; 50 ], seguida pelo intervalo ( 50 ; 60 ].
Destaca -se também, como um máximo local, o intervalo [ 0 ; 10 ], com distribuição em torno de 10 \% do total de notas, sendo que no intervalo [ 0 ; 10 ] inclui -se, além da nota zero, a frequência de alunos que deixaram esse tipo de questão em branco.
Na sequência, os resultados verificados para cada uma das questões discursivas de Formação Geral serão apresentados, estabelecendo -se relações com os conteúdos abordados em cada uma delas.
Os comentários da Banca de docentes corretores, a respeito do observado na correção das respostas dos estudantes, suas impressões e conclusões serão apresentados junto à análise de cada questão.
Cumpre esclarecer que, tendo -se em vista que as questões discursivas de Formação Geral são padronizadas, ou seja, constam de todas as provas, os comentários da Banca são os mesmos para todas as carreiras acadêmicas, sendo direcionados a todos os estudantes que participaram do Enade/2018.
A seguir, serão analisados os desempenhos da Área de Design nas duas questões discursivas de Formação Geral do Enade/2018, comparando -se os resultados obtidos com comentários para cada questão. 6.3.1.1 Análise de Conteúdo da Questão Discursiva 1 do Componente de Formação Geral Os dados de Design, obtidos a partir das respostas à questão 1, encontram -se na Tabela 6.14 e no Gráfico 6.7.
Nessa questão – a de melhor desempenho entre as duas de Formação Geral –, os alunos de todo o Brasil tiveram Média 39,6.
A maior Média para a questão 1 foi obtida na região Sudeste ( 40,7 ), e a menor, na região Norte ( 32,5 ).
Quanto à variabilidade das notas, o Desvio padrão de todo o Brasil foi 27,9.
O menor Desvio padrão foi obtido na região Centro-Oeste ( 27,4 ), e o maior Desvio padrão foi obtido na região Norte ( 29,3 ).
A Mediana do Brasil ( 40,0 ) foi igual a de três regiões : Nordeste, Sudeste e Centro-Oeste.
Nas outras duas regiões, as medianas foram menores : 30,0, na região Norte e 35,0, na região Sul.
As notas Máxima e Mínima da questão discursiva 1 foram, sem exceção, as mesmas para todas as regiões do Brasil, respectivamente, 100,0 e zero.
No Gráfico 6.7, mostra -se a distribuição das notas na questão discursiva 1 do Componente de Formação Geral.
A frequência modal foi o intervalo ( 30 ; 40 ], concentrando pouco mais de 15 \% dos participantes.
Caracteriza -se como máximo local a frequência que corresponde à dos alunos que deixaram a resposta a essa questão em branco ( 11,2 \% ), que somados àqueles que, tendo respondido à questão, obtiveram nota zero, ( 4,6 \% ) superam o valor modal. 6.3.1.2 Comentários sobre a correção de Conteúdo das respostas à Questão Discursiva 1 A Questão 1 exigia do aluno um posicionamento sobre conteúdo humanista e crítico.
O comando da questão solicitava a redação de um texto que formulasse e articulasse argumentos consistentes envolvendo temas previstos na Portaria nº 444, de 30 de maio de 2018, que dispõe sobre o componente de Formação geral do Enade 2018, tais como : solidariedade/violência, tolerância/intolerância, inclusão/exclusão, sexualidade, relações de gênero e relações \textbf{étnico-raciais} ( art. 7º, item VIII ), bem como ética, democracia e cidadania ( art. 7º, item I ).
O enunciado da questão buscava o estabelecimento da relação de três dimensões – os Direitos Humanos, a Democracia e as Fake News.
Para isso, recorreu a quatro suportes : a reprodução de uma foto e três fragmentos adaptados de textos que foram disponibilizados em páginas da Internet ( da Anistia Internacional, do Jornal El País-Brasil e do Labic-Laboratório de estudos sobre Imagem e Cibercultura ).
Os fragmentos traziam dados sobre o aumento do número de assassinatos de grupos minoritários, as origens e o trabalho da vereadora Marielle Franco, sua luta pelos Direitos Humanos, seu assassinato, e a produção de fake news.
O comando solicitava um texto abordando dois tópicos : o tensionamento entre a defesa dos Direitos Humanos e a produção de notícias falsas, tendo como base a vereadora Marielle – a defesa dos Direitos Humanos realizada por ela e as notícias falsas divulgadas após a sua morte.
As fake news também estavam presentes no segundo tópico de análise : os prejuízos que provocam para a sociedade democrática.
Por um lado, os textos selecionados para constituir o enunciado da questão continham apoio suficiente para formulação e articulação de argumentos para discorrer sobre o tensionamento entre a atuação da Marielle em defesa dos Direitos Humanos e fake news.
Por outro, havia ênfase nos aspectos biográficos e na trajetória da vereadora Marielle, o que levou a respostas pouco focadas no comando, que discutiam a vida e a morte da vereadora, sua importância como um ícone dos Direitos Humanos, opiniões em torno do seu assassinato, até mesmo sobre a responsabilidade pelo crime, muitos textos com a presença/defesa de fake news ou produção de novas.
O foco sobre Marielle também ensejou respostas basicamente \textbf{centradas} na questão da violência.
Foram frequentes as interpretações com viés político-partidário, não previstas no padrão de resposta.
O item do comando que solicitava a discussão sobre democracia e fake news encontrou menos apoio nos textos.
Apesar de se esperar que graduandos tivessem maturidade, inclusive política, para demonstrar alguma compreensão dos aspectos principais que envolvem o convívio democrático, esse aspecto foi pouco discutido e, quando presente, quase sempre foi apenas mencionado, numa quase reprodução do comando.
A falta de entendimento do que é o conceito de Direitos Humanos também foi percebida nas respostas, seja para recomendá -lo como importante para a sociedade, seja para criticá -lo como algo que defende bandidos.
Foi marcante a ideia que vincula Direitos Humanos a direitos apenas das minorias.
Não foi incomum encontrar respostas que \textbf{personificavam} os Direitos Humanos, tratando -os como sujeito de alguma ação que deveria ser implementada na sociedade em geral ou em casos específicos.
As fake news foram citadas muitas vezes, mas pouco explicadas, conforme previsto no padrão de resposta.
O argumento utilizado geralmente vinculava as consequências de fake news para a imagem da Marielle, ou seja, centrava no fato em si, já indicado no enunciado, vereadora Marielle Franco, sua luta pelos Direitos Humanos, seu assassinato, e a produção de fake news.
O comando solicitava um texto abordando dois tópicos : o tensionamento entre a defesa dos Direitos Humanos e a produção de notícias falsas, tendo como base a vereadora Marielle – a defesa dos Direitos Humanos realizada por ela e as notícias falsas divulgadas após a sua morte.
As fake news também estavam presentes no segundo tópico de análise : os prejuízos que provocam para a sociedade democrática.
Por um lado, os textos selecionados para constituir o enunciado da questão continham apoio suficiente para formulação e articulação de argumentos para discorrer sobre o tensionamento entre a atuação da Marielle em defesa dos Direitos Humanos e fake news.
Por outro, havia ênfase nos aspectos biográficos e na trajetória da vereadora Marielle, o que levou a respostas pouco focadas no comando, que discutiam a vida e a morte da vereadora, sua importância como um ícone dos Direitos Humanos, opiniões em torno do seu assassinato, até mesmo sobre a responsabilidade pelo crime, muitos textos com a presença/defesa de fake news ou produção de novas.
O foco sobre Marielle também ensejou respostas basicamente \textbf{centradas} na questão da violência.
Foram frequentes as interpretações com viés político-partidário, não previstas no padrão de resposta.
O item do comando que solicitava a discussão sobre democracia e fake news encontrou menos apoio nos textos.
Apesar de se esperar que graduandos tivessem maturidade, inclusive política, para demonstrar alguma compreensão dos aspectos principais que envolvem o convívio democrático, esse aspecto foi pouco discutido e, quando presente, quase sempre foi apenas mencionado, numa quase reprodução do comando.
A falta de entendimento do que é o conceito de Direitos Humanos também foi percebida nas respostas, seja para recomendá -lo como importante para a sociedade, seja para criticá -lo como algo que defende bandidos.
Foi marcante a ideia que vincula Direitos Humanos a direitos apenas das minorias.
Não foi incomum encontrar respostas que \textbf{personificavam} os Direitos Humanos, tratando -os como sujeito de alguma ação que deveria ser implementada na sociedade em geral ou em casos específicos.
As fake news foram citadas muitas vezes, mas pouco explicadas, conforme previsto no padrão de resposta.
O argumento utilizado geralmente vinculava as consequências de fake news para a imagem da Marielle, ou seja, centrava no fato em si, já indicado no enunciado, sem qualquer análise ou aprofundamento próprio.
É importante salientar que muitas respostas, provavelmente de grupos vinculados a cursos específicos, continham argumentações recorrentes : os Direitos Humanos como algo vinculado à Constituição e a tratados dos quais o Brasil é signatário ; as fake news como produto do mundo digital e globalizado, relacionado à era da pós-verdade ; a discussão sobre o papel das mídias sociais e da imprensa.
Como já apontado, houve baixa quantidade de respostas contemplando o segundo tópico do comando – prejuízos da produção de notícias falsas para a sociedade democrática ; fato motivado, possivelmente, pela falta de um elemento que ressaltasse esse ponto nos textos motivadores.
De modo geral, os estudantes não conseguiram recorrer a conhecimentos tácitos sobre democracia para construir argumentos próprios.
As respostas que abordavam o tema, geralmente estavam vinculadas à ideia do direito à verdade, do direito à integridade da pessoa ou dos direitos gerais garantidos pela Constituição.
A relação da questão com o momento político brasileiro foi indicada em um número expressivo de respostas.
Por vezes, apenas questionando a temática da questão – caso Marielle – em uma prova para avaliar o Ensino Superior no Brasil ou a importância dos Direitos Humanos na atual conjuntura política do país, levando à anulação da resposta.
Houve também interpretações das consequências do uso das fake news no processo eleitoral de 2018, influenciando a escolha consciente e livre das pessoas.
As respostas demonstraram atenção político social frente a esses temas, porém com pouco domínio conceitual do campo de conhecimento sobre Direitos Humanos, fake news e democracia.
Em termos da linguagem adotada pelos alunos, há dois grandes grupos de respostas.
No primeiro grupo, os que responderam de forma clara, com autonomia de pensamento e explicitação de justificativas para os aspectos solicitados no enunciado, demonstrando conhecimento sobre o que significam Direitos Humanos, fake news e democracia, além da articulação entre essas questões.
Por outro lado, um segundo grupo dos que não conseguiram ir além do que foi citado no enunciado, apenas repetindo -o, muitas vezes desviando o foco da questão para temas correlatos, como a vida de Marielle ou a violência.
Nesse segundo grupo, mais facilmente foram encontradas respostas desconexas e mal articuladas, com textos confusos que explicitam o não conhecimento sobre o tema e a falta de adequação à linguagem, necessária a um concluinte do Ensino Superior.
As abordagens dos Direitos Humanos mais recorrentes nas respostas corretas foram : a referência da Marielle como representante/ícone dos DH ; e os DH como direitos previstos na Constituição e em tratados internacionais dos quais o Brasil é signatário.
No entanto, também foram encontradas críticas aos Direitos Humanos, tais como : DH vistos como “ direitos de bandidos ” ; a defesa pouco embasada de frases de efeito como “ direitos humanos para humanos direitos ”.
Também chamou atenção o tratamento dos DH como um personagem, ou um órgão público, evidenciando desconhecimento sobre o tema.
Em relação às fake news, dois argumentos previstos pelo padrão de respostas foram os mais presentes : a noção de que o mundo globalizado e a facilidade de reprodução de informações pelas mídias atuais levam à criação de \textbf{mentiras} que são compartilhadas e que geram consequências ; a Marielle como alvo de fake news que tentaram desconstruir sua luta e seu legado.
Também foi frequente a ideia de que a produção de informações falsas é realizada por uma imprensa tendenciosa com intenção de manipular a informação.
No item relacionado a fake news e democracia, as respostas corretas afirmavam que as fake news inviabilizam o acesso à verdade e à garantia da integridade de um estado Democrático de Direito.
Também foi abordada a influência das fake news nos resultados das eleições dos últimos anos, nos estados Unidos e no Brasil, gerando consequências para a vida das pessoas.
Por fim, um outro aspecto abordado nas respostas a este tópico foi a polarização da sociedade em termos políticos, com a afirmação de que as fake news dificultam o debate, estimulam posições extremas e rivalidades prejudiciais à convivência em uma democracia plena.
A notas mais fracas ( de zero a 35 ) foram cerca de 20 \% das respostas corrigidas – as que não estavam em branco nem foram desconsideradas ou anuladas.
Foram respostas que se restringiam a cópia ou paráfrases do enunciado ; as que tinham foco em desdobramentos políticos e policiais do caso Marielle ; as que tinham argumentação confusa e/ou incompleta.
A grande maioria das notas, em torno de 65 \%, foi mediana, entre 40 a 70.
Nesse grupo ficaram as respostas que continham alguma abordagem adequada sobre Direitos Humanos, sobre as atividades da vereadora Marielle e sobre as fake news divulgadas após sua morte.
As notas mais altas deste grupo foram para respostas que incluíram a questão do tensionamento, quase sempre, em um texto coerente e consistente.
Muitas vezes, a nota mediana foi consequência do baixo índice de respostas à segunda parte do comando, que solicitava a discussão sobre os prejuízos das fake news para a democracia.
As melhores notas só foram obtidas por aqueles que incluíram a discussão sobre os prejuízos das fake news para a democracia.
Apenas 15 \% das respostas corrigidas receberam avaliação igual ou superior a 75.
A deficiência principal detectada pelos corretores foi a dificuldade de grande parte dos alunos para interpretar o enunciado da questão e para se expressar em textos consistentes.
De modo geral, os temas foram tratados de forma fragmentada e pouco analítica.
Além disso, chamou à atenção a predominância do lugar comum na construção argumentativa das respostas, a superficialidade do debate em torno da democracia e o desconhecimento de conceitos como Direitos Humanos, esses confundidos com instituições e até pessoas. 6.3.1.3 Análise de Conteúdo da Questão Discursiva 2 do Componente de Formação Geral Na Tabela 6.15 mostra -se que o desempenho médio dos estudantes, na questão discursiva 2 ( média 36,0 ), foi inferior ao obtido na questão discursiva 1 ( média 39,6 ).
A região Centro-Oeste foi aquela cuja Média, também nessa questão, foi maior ( 37,6 ), enquanto a de menor Média novamente ocorreu na região Norte ( 33,5 ).
Quanto à variabilidade das notas, o Desvio padrão de todo o Brasil foi 22,8, inferior ao obtido na questão discursiva 1 ( 27,9 ).
O maior desvio nessa questão foi obtido na região Centro-Oeste ( 24,2 ), enquanto o menor foi obtido na região Nordeste ( 22,4 ).
A Mediana de todo o Brasil foi 35,0, a mesma para todas as regiões.
A nota Máxima 100,0 foi obtida em três regiões : Sudeste, Sul e Centro-Oeste.
Na região Norte, a nota Máxima foi 85,0 e na região Nordeste, foi 90,0.
A nota Mínima foi zero em todas as regiões, sem exceção.
No Gráfico 6.8, mostra -se a distribuição das notas na questão discursiva 2 do Componente de Formação Geral.
A frequência modal ocorreu no intervalo ( 40 ; 50 ], com um quinto dos participantes, seguido do intervalo ( 30 ; 40 ].
Observa -se um máximo local na frequência que corresponde à dos alunos que deixaram a questão em branco, com mais de 11 \% dos estudantes.
Nota -se, ainda, que as notas ficaram menos dispersas em comparação às da questão discursiva de número 1, o que pode ser constatado também pela comparação do desvio padrão das notas da questão discursiva 2 ( 22,8 ) e o das notas da questão discursiva 1 ( 27,9 ). 6.3.1.4 Comentários sobre a correção de Conteúdo das respostas à Questão Discursiva 2 Os textos e o comando da Questão 2 de Formação Geral são perfeitamente adequados ao que se espera avaliar em Formação Geral, ao tratar das funções sociais dos museus, contextualizando a temática sob o ponto de vista da memória e da perspectiva de futuro.
Na questão, estavam presentes, principalmente, o conteúdo humanista e crítico, a interpretação e a produção de textos sobre temas como cultura e arte, educação, ciência e democracia, previstos na Portaria que dispõe sobre o componente de Formação Geral do Enade 2018.
Os conhecimentos exigidos para responder ao comando, explicitados no padrão de resposta, vão além do senso comum, por se exigir a compreensão de conceitos como “ função social ” ( de instituições sociais ), “ memória ” ( no sentido de memória social ), “ perspectiva de futuro ” ( no contexto \textbf{sócio-histórico-cultural} ) e “ sociedade contemporânea ” ( que deveria independer da interpretação de cada grupo social a respeito da sociedade na qual vive ).
No entanto, apesar de envolver tais conceitos, espera -se que alunos concluintes de cursos superiores de graduação sejam capazes de discutir funções de instituições sociais, especialmente porque o enunciado trazia fato recente e de grande repercussão como apoio.
O padrão de resposta elencou as funções sociais que poderiam ser abordadas nas respostas dos estudantes, como também argumentos que poderiam ser usados ao longo das mesmas.
A tarefa dos estudantes era a construção de um texto expositivo que articulasse a importância de museus, suas funções sociais e as perspectivas de memória e de futuro, o que se mostrou de dificuldade de média a difícil.
Apesar de o enunciado ser claro, muitos estudantes, provavelmente influenciados pelos textos, limitaram -se a discutir aspectos relacionados ao Museu Nacional e, em particular, ao incêndio que o \textbf{destruiu}.
Foram poucos os que interpretaram corretamente o escopo do comando, que solicitava a discussão das funções sociais de museus de uma forma geral, no sentido de argumentar a importância deste tipo de Instituição.
As quatro funções sociais elencadas no padrão de resposta foram encontradas na quase totalidade das redações dos estudantes de forma articulada com a importância dos museus.
Entretanto, quase nunca os textos abordavam mais do que uma função social e vinculavam, de forma clara, a importância desse tipo de Instituição para a sociedade contemporânea.
Além disso, a solicitação de contextualização sob o foco museu/memória/perspectiva de futuro parece não ter sido bem compreendida e, em consequência, esteve pouco presente nas respostas.
Também se observou confusão entre os conceitos de história e memória e uma visão estreita do significado de sociedade contemporânea, às vezes, restrita à cidade do respondente, outras, ao Rio de Janeiro, ou ainda a um tipo de sociedade idealizada.
A relevância que os textos de apoio deram ao caso particular do incêndio do Museu Nacional do Rio de Janeiro e aos prejuízos causados ao seu acervo, induziu a maioria dos estudantes a elaborar respostas tratando das funções a partir da perda e, em particular, à perda do tipo de acervo daquele museu e do tipo de pesquisa que lá se desenvolve, divergindo do comando da questão, que solicitava texto abordando funções sociais de museus em geral.
Também foram encontradas muitas respostas que se limitavam a discutir o que o estudante considerava como possíveis causas do incêndio, a falta de manutenção, questões políticas e de gestão, sem responder ao que foi solicitado.
Foi também recorrente a transcrição de frases dos textos do enunciado nas respostas, diretamente ou por paráfrases, revelando certa dificuldade de criação textual.
A análise das respostas dos estudantes, além de revelar dificuldade de abordar a temática esperada, revelou outros dois aspectos predominantes.
Um deles, como já comentado, foi a ênfase no incêndio do Museu Nacional, que orientou grande parte dos textos dos estudantes.
Outro, também evidenciando desvio do esperado, foi o foco na perda patrimonial.
Nesse grupo, houve respostas que até apresentaram uma função social, mas olhada pelo ponto de vista da perda.
Nessa linha, muitas respostas enalteciam a importância do Museu Nacional, sua relevância no cenário nacional, o impacto da perda de seu acervo, isso complementado por frágeis reflexões sobre as funções sociais.
Em consequência, a perspectiva de desdobramentos mais ricos, que tratassem de aspectos sociais relacionados a diferentes tipos de museu foi prejudicada.
Muitas respostas trataram da preservação da memória e/ou guarda do acervo ; identidade, tanto individual quanto coletiva ; conhecimento sobre antepassados ; evolução da sociedade e dos seus costumes ; valorização e disseminação do conhecimento por meio da liberação do acesso aos acervos ; bem como atividades pedagógicas e exposições de arte.
O aspecto produção de pesquisa e conhecimento, embora frequente, foi mencionado de forma muito simples, sem qualquer aprofundamento, para além do que já era mencionado no enunciado.
Frequentemente apresentava -se o museu como produtor de pesquisas para o campo da saúde, o que evidencia desconhecimento das reais possibilidades de pesquisa, mesmo no Museu Nacional.
Quanto a perspectiva de futuro, quando presente no texto, era mencionada de forma breve, com afirmações genéricas sobre a possibilidade de melhoria do futuro pelo simples conhecimento do passado, numa percepção de que o passado determina o futuro, de forma linear.
Foram raras as respostas com alguma argumentação que associasse pesquisa, memória, identidade e reflexão social como base para a construção de alternativas de futuro.
A relação entre turismo e museu apareceu considerando a visita a museus como forma de lazer e, algumas vezes, como geradora de recursos.
O desempenho pode ser considerado fraco.
Muitas respostas deixavam claro que o estudante nunca tinha visitado um museu e, de uma forma geral, observou -se desconhecimento sobre suas funções sociais.
Não se esperava que história, memória social e museologia fossem assuntos dominados pelos estudantes das diversas áreas que participaram do Enade 2018, mas esperava -se que essa temática pudesse ser abordada do ponto de vista das experiências pessoais ou de conhecimento social.
Como uma vivência própria com os museus também se revelou frágil ou inexistente, o conhecimento sobre o assunto ficou bem restrito aos textos motivadores e ao senso comum, construído, basicamente, pelo noticiário em torno do incêndio.
O universo dos argumentos dos estudantes se mostrou muito limitado.
Cada dimensão foi praticamente representada por uma ideia simplista.
De forma recorrente : a preservação da memória para construção da identidade de um povo ; a compreensão do passado por meio das pesquisas como oportunidade para melhorar o futuro ; o museu como um espaço em que há um \textbf{aprendizado} estimulado pelo lazer. das pesquisas como oportunidade para melhorar o futuro ; o museu como um espaço em que há um \textbf{aprendizado} estimulado pelo lazer.
Houve, também, uma grande concentração de respostas associando o museu a um local de guarda de acervo e de história.
As respostas à Questão 2 evidenciaram dificuldades conceituais.
Além dos conceitos de história, memória e museu se misturarem, parecendo não se distinguirem, ao mencionarem a pesquisa, também se observou falta de distinção entre pesquisas colegiais, com objetivos de \textbf{aprendizado}, e as realizadas por pesquisadores, com objetivos de avançar cientificamente.
A falta de compreensão do que se pedia no comando da questão e a dificuldade de expressão linguística foram dois aspectos marcantes na avaliação dos corretores.
Foi frequente os textos apresentarem ideias soltas ou de forma confusa, sem constituir unidade.
O vocabulário usado se revelou bem restrito, e fortemente apoiado no coloquial, na linguagem oral.
Predominou o senso comum em lugar de uma reflexão mais formal envolvendo os conceitos envolvidos. 6.3.1.5 Análise de Língua Portuguesa das Questões Discursivas do Componente de Formação Geral Os dados de Design, obtidos a partir das respostas às questões discursivas do Componente de Formação Geral, no que tange à Língua Portuguesa, encontram -se na Tabela 6.16 e no Gráfico 6.9.
Nesse aspecto, os alunos de todo o Brasil obtiveram Média 55,8.
A maior Média em relação à Língua Portuguesa foi obtida na região Centro-Oeste ( 57,6 ), e a menor, na região Norte ( 49,6 ).
Quanto à variabilidade das notas, o Desvio padrão de todo o Brasil foi 22,3.
O menor Desvio padrão foi obtido na região Sudeste ( 21,2 ) e o maior Desvio padrão foi obtido na região Norte ( 25,6 ).
A Mediana das notas de Língua Portuguesa foi 62,5 para o Brasil, como um todo, a mesma obtida nas regiões Sudeste e Sul.
A Mediana das regiões Norte e Nordeste foi 60,0, enquanto a da Centro-Oeste foi 65,0, a maior dentre as obtidas pelas Grandes Regiões.
A nota Máxima para todo o Brasil foi 95,0, com, pelo menos, um aluno tirando essa nota na região Nordeste.
Nas regiões Norte e Sul, a nota Máxima foi 92,5 ; na Sudeste, foi 90,0 ; e, na região Centro-Oeste, foi 85,0, a menor nota Máxima dentre as regiões.
Já a nota Mínima foi zero em todas as regiões do país.
No Gráfico 6.9, mostra -se a distribuição das notas de Língua Portuguesa do Componente de Formação Geral.
Observa -se que a maior frequência ( quase 30 \% ) corresponde à dos alunos que obtiveram nota no intervalo ( 60 ; 70 ].
Destacam -se, também, os alunos que deixaram ambas as questões em branco, representando 7,5 \% do total, caracterizando -se como máximo local. 6.3.1.6 Comentários sobre a correção das respostas de Formação Geral com respeito à Língua Portuguesa As duas questões discursivas de Formação Geral ensejaram duas configurações textuais distintas, tendo em vista os seus encaminhamentos : a Questão 1 teve um caráter mais opinativo e a Questão 2, mais expositivo, conforme se explicita a seguir.
Ambas as questões permitiram a \textbf{observação} do desempenho linguístico dos participantes em situação formal, competência exigida no processo de formação universitária.
As habilidades de produção textual envolvidas têm, igualmente, como condição básica, o domínio do padrão formal culto da Língua Portuguesa, no que diz respeito aos aspectos relativos às questões de caráter ortográfico, textual, morfossintático e vocabular.
O encaminhamento das questões evidencia a expectativa de que o participante utilizasse seus conhecimentos sobre os dois temas e estruturasse seus textos de acordo com as características do registro formal adequado à situação comunicativa.
Tal configuração determina exigências quanto : à adequação da seleção vocabular, ao desenvolvimento do conteúdo, à estruturação sintática dos períodos, à organização lógica das ideias, à utilização de procedimentos de encadeamento textual e de referenciação, à obediência às exigências morfossintáticas próprias da modalidade escrita da norma-padrão, ao respeito às regras ortográficas e às regras de acentuação gráfica.
Esse encaminhamento exigiu do graduando o domínio do modo de organização textual de base dissertativo-argumentativa, essencial ao processo de formação universitária.
As habilidades de produção textual envolvidas têm como condição básica o domínio do padrão formal culto da Língua Portuguesa, no que diz respeito aos aspectos relativos às questões de caráter ortográfico, textual, morfossintático e vocabular.
O enunciado da Questão 1 tomava como base três fragmentos de textos jornalísticos publicados em sites da internet, que abordavam o aumento do assassinato de minorias e de defensores de direitos humanos no Brasil, entre eles a vereadora Marielle Franco.
O comando solicitava que o participante desenvolvesse um texto dissertativo e argumentativo com as temáticas Fake News, Democracia e Direitos Humanos.
Embora não tenha havido referência explícita à tipologia textual esperada para o desenvolvimento da resposta, essa proposta encaminhou o participante para a elaboração de um texto opinativo, com o objetivo de discutir o problema e manifestar um posicionamento crítico.
O enunciado da Questão 2 abordava o incêndio do Museu Nacional do Rio de Janeiro, ocorrido em 2018, e seus impactos na ciência brasileira e internacional, tomando como base dois fragmentos de textos jornalísticos publicados em sites da internet.
O comando solicitava que o participante desenvolvesse um texto abordando três aspectos da função social dos museus na sociedade contemporânea, encaminhando -o para a elaboração de um texto expositivo que confirmasse a posição defendida nos fragmentos motivadores a respeito da importância dessas instituições.
Para atender a essas exigências, o padrão de resposta utilizado na avaliação das questões discursivas de Formação Geral considerou aspectos relevantes ao bom desempenho linguístico como competências distintas, de modo a permitir um mapeamento detalhado do domínio dos recursos disponíveis na Língua Portuguesa para a comunicação escrita formal : ortográficos, textuais, morfossintáticos e vocabulares ( ver detalhes no Anexo VIII – Padrão de Respostas ).
Aspectos Ortográficos Para avaliar a competência relativa ao domínio das convenções ortográficas da norma-padrão da Língua Portuguesa, observou -se o respeito às regras de acentuação gráfica e da grafia padrão das palavras ( com ausência de abreviaturas próprias da linguagem da internet ), de acordo com as convenções estabelecidas pela legislação em vigor e consubstanciadas no Vocabulário Ortográfico da Língua Portuguesa, editado pela Academia Brasileira de Letras.
Aspectos Textuais Para avaliar a competência relativa ao domínio dos procedimentos de estruturação textual, procurou -se observar a utilização de mecanismos linguístico-discursivos que estabeleçam o encadeamento lógico entre as partes do texto, de forma a garantir a progressão e a coerência textuais.
Esta competência destaca -se por sua relevância, já que são avaliadas a clareza e a unidade textual.
Aspectos Morfossintáticos Para avaliar o domínio dos diferentes aspectos morfossintáticos próprios da modalidade escrita formal da norma-padrão da Língua Portuguesa, observou -se se o texto produzido atendeu às seguintes exigências : a concordância nominal, a concordância verbal, a regência nominal, a regência verbal, a flexão nominal, a flexão verbal, a correlação entre os tempos verbais, a colocação pronominal e a utilização de sinais de pontuação que contribuam para a organização lógica da frase e do texto.
Aspectos Vocabulares Para avaliar a adequação da seleção vocabular à modalidade escrita formal da Língua Portuguesa, exigida pela situação comunicativa, observou -se o atendimento às seguintes exigências : precisão na seleção/utilização do vocabulário relacionado à temática solicitada pela questão ; ausência de marcas da oralidade, como termos de sentido muito genérico ( “ coisa ”, “ negócio ”, “ você ” ) e termos de registros mais informais ( como gírias, jargões, frases feitas, ditados populares, termos regionais ).
Assim, espera -se que o participante respeite a adequação vocabular não usando gírias ou expressões coloquiais, evite repetição desnecessária de palavras e utilize um vocabulário mais formal, como solicitado por um texto dissertativo.
Para efeito de pontuação, a avaliação do desempenho linguístico considerou três grandes grupos de competências, agrupando aspectos explicitados, com os seguintes \textbf{pesos} relativos : aspectos ortográficos ( 20 \% ) ; aspectos textuais ( 40 \% ), aspectos morfossintáticos e vocabulares ( 40 \% ).
Os critérios de pontuação utilizados no padrão de resposta procuraram atender à diversidade de desempenho dos graduandos.
Nos textos analisados, observaram -se, assim, vários níveis de atendimento às exigências do padrão formal próprio da modalidade escrita da Língua Portuguesa, configurando um"continuum".
Por um lado, textos que traduzem com exatidão a expectativa de uso formal da língua, em todos os níveis analisados, marcados por clareza e unidade ; e, no outro extremo, textos quase incompreensíveis devido ao grande número de desvios de todos os tipos, principalmente os de caráter textual.
As duas questões discursivas de Formação Geral analisadas revelaram desempenhos distintos dos participantes.
A Questão 1 propiciou a oportunidade de maior desenvolvimento do tema solicitado, gerando consequências na estruturação textual, já que os textos foram mais longos, mais elaborados e fluentes.
Em virtude da polêmica sobre a existência das fake news e sobre a visão pública da personagem da vereadora Marielle Franco, observaram -se respostas que polemizaram o tema, favorecendo a construção argumentativa do texto.
A Questão 2 teve um comportamento distinto, devido ao caráter mais expositivo do enunciado e ao reduzido conhecimento dos participantes sobre o tema.
Assim, os textos foram mais curtos, frequentemente com repetição de ideias dos textos motivadores.
Quanto aos aspectos linguísticos analisados durante esta avaliação, observaram -se os seguintes resultados : Aspectos ortográficos : O desempenho dos participantes revelou uma diferença muito grande nos dois aspectos analisados nesta competência : baixo índice de desvios da grafia padrão e grande índice de desvios de acentuação.
Em vários casos, ocorre ausência completa de acentuação gráfica em todas as palavras do texto.
Os resultados revelaram que a tendência dominante entre os universitários brasileiros é a eliminação da acentuação gráfica, talvez motivada pelos hábitos relacionados às redes sociais e pela ausência de esclarecimento dos meios de comunicação, das autoridades e das escolas sobre as decisões do Acordo Ortográfico da Língua Portuguesa de 1990.
Os casos mais sistemáticos de eliminação do acento indicador da sílaba tônica foram : - palavras proparoxítonas ( por exemplo “ generos ”, “ arbitro ”, “ politico ”, “ publicas ”, “ numero ” ) ; - palavras paroxítonas terminadas em ditongo crescente ( por exemplo “ area ”, “ noticiarios ”, “ individuo ”, “ historia ”, “ varios ”, “ noticia ”, “ negligencia ”, “ ocorrencia ” ) ; - palavras oxítonas ( por exemplo “ ninguem ”, “ esta ”, “ ate ”, “ tambem ” ).
Por outro lado, destaca -se o uso indevido do acento gráfico em palavras como “ melâncolia ” e “ intervênção ”, por exemplo.
Quanto ao domínio das convenções relativas à grafia das palavras, observaram -se desvios \textbf{esporádicos} como : hipercorreção pela escolha de “ e ” no lugar de “ i ”, por influência de hábitos da oralidade ( “ descriminação ” no lugar de “ discriminação ”, “ entervenção ” no lugar de “ intervenção ” ) ; desvios de grafia relacionados à variação diastrática, como por exemplo “ perca ”, “ subjulgadas ”, “ precoseito ”, “ soubre ”, “ vecendo ”, “ indesencia ”, “ apolojia ”, “ fulga ”, “ dereitos ” ;...
Observaram -se, também, muitos casos de inadequação no uso da maiúscula : uso indevido para destacar determinadas palavras-chave do texto, como “ Brasileiros ”, “ Homens ” ; ausência de maiúscula para grafar nomes próprios, como"marielle".
Expressões como"Direitos Humanos"foram grafadas diversificadamente, com maiúsculas e com minúsculas.
Destaque -se, também, a maioria dos participantes que grafam os textos inteiramente em caixa alta, o que impede a identificação de desvios relacionados ao uso de maiúsculas.
Vale observar, também, que, ao contrário do que se esperava, não apareceram abreviaturas próprias do “ internetês ”, relacionadas ao uso de redes sociais e e-mails, como por exemplo"vc".
Aspectos textuais : Para alguns dos graduandos essa competência se revelou como a mais problemática, tendo em vista os inúmeros problemas observados, desvios acumulados durante toda a formação escolar.
São eles : sequência justaposta de ideias sem encaixamentos sintáticos, provocando uma redução de estruturas subordinadas, ao lado do aumento na frequência de estruturas coordenadas e absolutas ; frases fragmentadas que comprometem a estrutura lógico-gramatical ; frases formadas apenas por oração subordinada, sem oração principal ; redução no uso de conectores para expressar relações lógicas essenciais à construção do texto como consequência da mudança de estruturação frasal ; emprego inadequado do pronome relativo ( com omissão da preposição ou a utilização do pronome relativo inadequado, como “ onde ” ), refletindo hábitos da oralidade ; ausência de recursos de referenciação, como substituição de termos por sinônimos, hiperônimos, nominalizações, expressões metafóricas.
Esses problemas revelam dificuldades em relação à estrutura formal do texto produzido, o que é preocupante ao se levar em conta que são graduandos em fase final de formação.
Em alguns textos, falta um mínimo de textualidade e de domínio do registro padrão da língua.
Na verdade, observaram -se relações linguísticas quase agramaticais, como as estabelecidas pela sequência de gerúndios sem o apoio de um ponto de partida para a organização das informações gramaticais e semânticas.
Quanto à utilização dos sinais de pontuação, observou -se uma ausência dos mesmos na maioria dos textos analisados.
Foi muito frequente a ocorrência de parágrafos sem marca interna de pontuação para separar os períodos.
São os seguintes os tipos de problemas mais encontrados :  vírgula : utilização de vírgula para separar o sujeito e o predicado ; uso de vírgula no lugar do ponto para separar ideias que constituem períodos distintos ; ausência de vírgula para separar oração adjetiva explicativa ;  ponto e vírgula : utilização do ponto e vírgula no lugar de vírgula ;  ponto final : ausência de ponto final para separar períodos.
Aspectos morfossintáticos e vocabulares : Os resultados são transparentes em relação aos aspectos mais problemáticos no desempenho dos participantes nestas duas competências.
Em relação aos aspectos morfossintáticos, seguem algumas \textbf{observações} sobre os desvios mais frequentes.
Quanto à regência, o desvio mais frequente foi a falta do sinal indicativo da crase – isso revela que o usuário não tem consciência de que, sob a forma do termo “ a ”, existe a presença de uma contração entre a preposição “ a ” ( exigida pela regência do termo anterior ) e o artigo definido “ a ”.
Outro problema relacionado à regência verbal e à nominal, encontrado frequentemente nas respostas, foi a ausência de preposição antes de pronome relativo, processo generalizado lógico-gramatical ; frases formadas apenas por oração subordinada, sem oração principal ; redução no uso de conectores para expressar relações lógicas essenciais à construção do texto como consequência da mudança de estruturação frasal ; emprego inadequado do pronome relativo ( com omissão da preposição ou a utilização do pronome relativo inadequado, como “ onde ” ), refletindo hábitos da oralidade ; ausência de recursos de referenciação, como substituição de termos por sinônimos, hiperônimos, nominalizações, expressões metafóricas.
Esses problemas revelam dificuldades em relação à estrutura formal do texto produzido, o que é preocupante ao se levar em conta que são graduandos em fase final de formação.
Em alguns textos, falta um mínimo de textualidade e de domínio do registro padrão da língua.
Na verdade, observaram -se relações linguísticas quase agramaticais, como as estabelecidas pela sequência de gerúndios sem o apoio de um ponto de partida para a organização das informações gramaticais e semânticas.
Quanto à utilização dos sinais de pontuação, observou -se uma ausência dos mesmos na maioria dos textos analisados.
Foi muito frequente a ocorrência de parágrafos sem marca interna de pontuação para separar os períodos.
São os seguintes os tipos de problemas mais encontrados :  vírgula : utilização de vírgula para separar o sujeito e o predicado ; uso de vírgula no lugar do ponto para separar ideias que constituem períodos distintos ; ausência de vírgula para separar oração adjetiva explicativa ;  ponto e vírgula : utilização do ponto e vírgula no lugar de vírgula ;  ponto final : ausência de ponto final para separar períodos.
Aspectos morfossintáticos e vocabulares : Os resultados são transparentes em relação aos aspectos mais problemáticos no desempenho dos participantes nestas duas competências.
Em relação aos aspectos morfossintáticos, seguem algumas \textbf{observações} sobre os desvios mais frequentes.
Quanto à regência, o desvio mais frequente foi a falta do sinal indicativo da crase – isso revela que o usuário não tem consciência de que, sob a forma do termo “ a ”, existe a presença de uma contração entre a preposição “ a ” ( exigida pela regência do termo anterior ) e o artigo definido “ a ”.
Outro problema relacionado à regência verbal e à nominal, encontrado frequentemente nas respostas, foi a ausência de preposição antes de pronome relativo, processo generalizado na modalidade oral da língua, em situações de registro informal.
Apesar da possibilidade de que essa alteração de regência se generalize no padrão escrito da Língua Portuguesa, como já está ocorrendo até em textos jornalísticos, o não emprego da preposição foi considerado como desvio neste processo de avaliação.
Outro desvio de regência significativo foi a utilização inadequada de uma preposição ou sua ausência após o verbo ou o nome ( substantivo ou adjetivo ).
A concordância verbal e a concordância nominal apresentaram alguns desvios muito frequentes.
Quanto à concordância de número, observou -se, como apontam as pesquisas nessa área, ausência de marca ( com sujeito anteposto ou posposto ) ou uso indevido ( uso inadequado da marca de plural comandado pelo núcleo plural da locução adjetiva, apesar de o substantivo que funciona como núcleo do sintagma nominal estar no singular ).
Uma ocorrência que se destacou foi a ausência de acento circunflexo na forma plural do presente do indicativo do verbo “ ter ”, que foi considerada como um desvio na concordância verbal e não na acentuação gráfica.
Quanto à concordância de gênero, vários casos foram observados, normalmente no âmbito de sintagmas nominais longos, em que o adjetivo está afastado do substantivo.
Quanto à questão da colocação pronominal, foram poucos os desvios observados.
Concluiu -se que, no registro escrito formal, a maioria dos participantes já incorporou regras como a não introdução da frase por um pronome oblíquo e a próclise na presença de um termo atrator.
Várias marcas de oralidade foram identificadas, embora não com alta frequência : uso do pronome relativo “ onde ” como relativo universal, falta de artigo definido antes de substantivo isolado ou antes de termos em enumerações.
Quanto aos aspectos vocabulares, alguns tipos de inadequação foram observados : expressões da oralidade ; seleção vocabular incompatível com o contexto, gerando falta de inteligibilidade ; falta de domínio de vocabulário mais abstrato e de maior complexidade, essencial ao desenvolvimento do texto de base dissertativa.
O principal aspecto observado foi a excessiva repetição de certas palavras, como o termo"pessoa", por exemplo, revelando limitação de repertório vocabular.
Em função do tema solicitado na Questão 1, diferentes grafias da expressão fake news não foram consideradas.
Para sintetizar, os problemas que mais se destacaram nos textos analisados foram : a ) truncamentos sintáticos, caracterizados por estruturas frasais incompreensíveis devido à ruptura da complexidade sintática própria da modalidade escrita ; na modalidade oral da língua, em situações de registro informal.
Apesar da possibilidade de que essa alteração de regência se generalize no padrão escrito da Língua Portuguesa, como já está ocorrendo até em textos jornalísticos, o não emprego da preposição foi considerado como desvio neste processo de avaliação.
Outro desvio de regência significativo foi a utilização inadequada de uma preposição ou sua ausência após o verbo ou o nome ( substantivo ou adjetivo ).
A concordância verbal e a concordância nominal apresentaram alguns desvios muito frequentes.
Quanto à concordância de número, observou -se, como apontam as pesquisas nessa área, ausência de marca ( com sujeito anteposto ou posposto ) ou uso indevido ( uso inadequado da marca de plural comandado pelo núcleo plural da locução adjetiva, apesar de o substantivo que funciona como núcleo do sintagma nominal estar no singular ).
Uma ocorrência que se destacou foi a ausência de acento circunflexo na forma plural do presente do indicativo do verbo “ ter ”, que foi considerada como um desvio na concordância verbal e não na acentuação gráfica.
Quanto à concordância de gênero, vários casos foram observados, normalmente no âmbito de sintagmas nominais longos, em que o adjetivo está afastado do substantivo.
Quanto à questão da colocação pronominal, foram poucos os desvios observados.
Concluiu -se que, no registro escrito formal, a maioria dos participantes já incorporou regras como a não introdução da frase por um pronome oblíquo e a próclise na presença de um termo atrator.
Várias marcas de oralidade foram identificadas, embora não com alta frequência : uso do pronome relativo “ onde ” como relativo universal, falta de artigo definido antes de substantivo isolado ou antes de termos em enumerações.
Quanto aos aspectos vocabulares, alguns tipos de inadequação foram observados : expressões da oralidade ; seleção vocabular incompatível com o contexto, gerando falta de inteligibilidade ; falta de domínio de vocabulário mais abstrato e de maior complexidade, essencial ao desenvolvimento do texto de base dissertativa.
O principal aspecto observado foi a excessiva repetição de certas palavras, como o termo"pessoa", por exemplo, revelando limitação de repertório vocabular.
Em função do tema solicitado na Questão 1, diferentes grafias da expressão fake news não foram consideradas.
Para sintetizar, os problemas que mais se destacaram nos textos analisados foram : a ) truncamentos sintáticos, caracterizados por estruturas frasais incompreensíveis devido à ruptura da complexidade sintática própria da modalidade escrita ; b ) redução de estruturas subordinadas, compensada pelo aumento na frequência de estruturas coordenadas e absolutas, por um lado, ou pela elaboração de estruturas textuais prejudicadas pelo excesso de ideias sem a devida conexão subordinativa ; c ) redução no uso de conectores para expressar relações lógicas essenciais à construção do texto, substituídos pela exigência de inferência por parte do interlocutor para suprir a sua ausência ; d ) redução no uso do subjuntivo, ao lado da ampliação do uso do indicativo, combinado a estruturas frasais coordenadas ou absolutas ; d ) empobrecimento do processo de referenciação, com a repetição de palavras, sem processos mais sofisticados de substituição ; e ) simplificação da marcação da categoria tempo na morfologia verbal ; f ) desrespeito às regras de concordância verbal e nominal próprias do padrão formal ; g ) ausência do acento grave indicador da crase, revelando falta de conhecimento dos critérios de regência verbal e nominal ; h ) falta de domínio de vocabulário mais abstrato e de maior complexidade, essencial ao desenvolvimento do processo dissertativo ; i ) redução do emprego da acentuação gráfica.
Como resultado da correção, merece atenção a constatação de que a modalidade escrita tem apresentado a tendência a uma evidente simplificação, aproximando -se das características da modalidade oral da Língua Portuguesa.
No caso dos textos de base dissertativa ( expositiva ) e de base argumentativa, inscritos em um registro formal, a distância entre as duas modalidades é ainda maior, o que provoca desvios recorrentes em todos os quatro aspectos analisados : ortográficos, textuais, morfossintáticos e vocabulares.
Destaca -se, nessa avaliação, o comprometimento textual, por meio de estruturas fragmentadas e/ou truncadas, rompendo a complexidade sintática esperada no padrão formal. das notas, o Desvio padrão de todo o Brasil foi 18,0.
O maior Desvio padrão foi encontrado na região Centro-Oeste ( 21,8 ), e o menor, na região Norte ( 15,0 ).
A maior nota Máxima, 96,7, foi obtida nas regiões Nordeste e Sudeste.
Nas demais regiões, as notas Máximas foram : 68,3, na região Norte e 91,7, nas regiões Sul e Centro-Oeste.
A nota Mínima ( 0,0 ) foi obtida por, pelo menos, um aluno em todas as regiões do Brasil.
A Mediana do Brasil, como um todo, foi 26,7.
A maior Mediana foi obtida na região Centro-Oeste ( 35,0 ), e a menor, na região Norte ( 21,7 ).
No Gráfico 6.10, representa -se a distribuição das notas nas questões discursivas no Componente de Conhecimento Específico.
A moda dessa distribuição ocorre no intervalo ( 20 ; 30 ], com mais de um quarto do total de participantes.
O intervalo [ 0 ; 10 ], onde estão computadas as ocorrências de respostas em branco para as três questões discursivas de Conhecimentos Específicos, chama a atenção com um total de 13,5 \% de estudantes.
A análise de cada uma destas questões será feita a seguir. 6.3.2.1 Análise da Questão Discursiva 3 do Componente de Conhecimento Específico Na questão 3, cujos resultados aferidos encontram -se descritos na Tabela 6.18, a Média dos estudantes de todo o Brasil foi 12,8, a questão de pior desempenho dentre as três discursivas de Conhecimentos Específicos.
A menor Média nessa questão foi obtida pelos alunos da região Norte ( 4,5 ), enquanto a maior Média foi obtida na região Centro-Oeste ( 15,4 ).
Quanto à variabilidade das notas, o Desvio padrão de todo o Brasil foi 21,4.
O maior Desvio padrão foi obtido na região Sudeste ( 22,7 ), enquanto o menor foi obtido na região Norte ( 12,2 ).
A nota Máxima, 100,0 pontos, foi alcançada por, pelo menos, um aluno de quase todas as regiões, com exceção da região Norte que teve nota Máxima 60,0.
A Mediana do Brasil, como um todo, foi zero, a mesma para todas as regiões.
A nota Mínima ( 0,0 ) também foi a mesma em todas as regiões do Brasil.
Cabe salientar que, pelo menos, a metade dos alunos em cada uma das regiões tirou nota zero nessa questão.
No Gráfico 6.11, mostra -se a distribuição das notas na questão discursiva 3, do Componente de Conhecimento Específico da área de Design.
A moda dessa questão foi a nota zero, para 50,8 \% dos participantes.
Somados aos que deixaram a questão em branco, chega -se a dois terços dos estudantes.
Sem levar em consideração esse contingente, a moda dessa distribuição seria o intervalo ( 10 ; 20 ], seguido do intervalo ( 20 ; 30 ]. 6.3.2.2 - Comentários sobre as respostas à Questão Discursiva 3 A questão discursiva 3 da prova de Design apresentava duas figuras de cartazes da III e da IV Bienal do MASP ocorridas, respectivamente, nos anos de 1955 e 1957.
O comando No Gráfico 6.11, mostra -se a distribuição das notas na questão discursiva 3, do Componente de Conhecimento Específico da área de Design.
A moda dessa questão foi a nota zero, para 50,8 \% dos participantes.
Somados aos que deixaram a questão em branco, chega -se a dois terços dos estudantes.
Sem levar em consideração esse contingente, a moda dessa distribuição seria o intervalo ( 10 ; 20 ], seguido do intervalo ( 20 ; 30 ]. 6.3.2.2 - Comentários sobre as respostas à Questão Discursiva 3 A questão discursiva 3 da prova de Design apresentava duas figuras de cartazes da III e da IV Bienal do MASP ocorridas, respectivamente, nos anos de 1955 e 1957.
O comando solicitava que o estudante citasse e descrevesse duas características do estilo gráfico e da linguagem desenvolvidos por Wollner, autor dos cartazes, ao longo de sua carreira.
O Enunciado estava claro, com linguagem apropriada e pertinente ao assunto tratado, no que se refere ao campo do design.
O conteúdo abordado foi adequado em relação às Diretrizes, pois concerne à História do Design Brasileiro, corpo de conhecimento fundamental para a formação do designer brasileiro.
A indicação para que o estudante descrevesse as características citadas confere à questão um nível maior de profundidade.
O enunciado possibilitava a ocorrência de interpretações variadas.
A reprodução de cartazes de autoria do designer citado contribuiu para a clareza e a objetividade do enunciado.
No entanto, levou a respostas com descrições exclusivamente formais e não de características do estilo gráfico e linguagem nos termos definidos no comando da questão.
A questão foi de média dificuldade em relação às demais questões.
Nas respostas, não houve muita diversidade na citação de estilo gráfico e linguagem.
As ocorrências verificadas de citação estavam previstas no padrão de resposta.
A única situação observada que poderia ser considerada como não prevista no padrão de resposta foi a menção a Bauhaus, presente em algumas respostas, o que não é incorreto.
O que surgiu com mais frequência nas respostas foi a citação de características que podiam ser facilmente visualizadas e identificadas nos exemplos apresentados, como o uso de formas geométricas, tipografia sem serifa, paleta de cor, minimalismo e abstração, sem descrições que evidenciassem o estilo gráfico e a linguagem desenvolvidos pelo designer Alexandre Wollner.
As respostas medianas fizeram, de forma superficial, menção ao Modernismo, ao Funcionalismo, a Bauhaus, referências que influenciaram o trabalho do designer citado.
O resultado final demonstra um alto índice de respostas fracas.
Verificou -se que a abordagem recorrente para responder à questão foi partir da descrição da imagem dos cartazes apresentados no enunciado.
Os principais erros se concentraram em descrever o que estava sendo visto na imagem do cartaz e não fazer relação com o estilo gráfico do designer.
Não houve diversidade relevante nas respostas.
As respostas consideradas fracas ( notas de 0 a 30 ) foram aquelas que apenas citaram as características e não descreveram o estilo ou ficaram apenas com a descrição da imagem do cartaz ( 71,1 \% ).
As respostas consideradas medianas ( notas de 35 a 70 ) foram as que citaram as características e fizeram uma descrição mediana, demonstrando um vocabulário incipiente para descrever o estilo gráfico do designer ( 11,0 \% ).
As respostas consideradas boas ( notas de 75 a 100 ) foram aquelas que citaram as características e descreveram adequadamente o estilo gráfico do designer, utilizando o vocabulário correto e, na grande maioria, utilizando palavras que constavam no padrão de resposta ( 2,4 \% ).
Um total de 15,5 \% de estudantes deixou a questão em branco. maioria, utilizando palavras que constavam no padrão de resposta ( 2,4 \% ).
Um total de 15,5 \% de estudantes deixou a questão em branco.
Em relação ao conteúdo exigido, o desempenho dos respondentes pode ser considerado insuficiente a fraco.
Identifica -se a tendência de descrição das formas que compõem a imagem do cartaz apresentado, a qual pode ter distraído os respondentes no foco central do enunciado.
Pode -se supor que, no ensino de História do Design, não há uma ênfase nos conteúdos relativos à História do Design Brasileiro. 6.3.2.3 - Análise da Questão Discursiva 4 do Componente de Conhecimento Específico Na Tabela 6.19, constam as informações relativas à questão 4 do conjunto de questões do Componente de Conhecimento Específico.
O desempenho dos estudantes nessa questão foi o melhor dentre as 3 discursivas.
A Média geral do Brasil foi 55,9, sendo a maior Média registrada na região Centro-Oeste ( 61,2 ), e a menor, na região Norte ( 49,6 ).
A nota Máxima ( 100,0 ) foi atingida por, pelo menos, um concluinte das regiões Nordeste, Sudeste e Sul.
No Norte e no Centro-Oeste, a nota Máxima foi 90,0.
A Mediana em todo o Brasil foi 60,0, o mesmo valor foi obtido na região Sudeste.
Nas demais regiões, a Mediana foi : 55,0, na região Norte ; 65,0, nas regiões Nordeste e Sul ; e 75,0, na região Centro-Oeste.
A nota Mínima foi zero em todas as regiões, sem exceção.
No Gráfico 6.12, está representada a distribuição de notas da questão discursiva 4, do Componente de Conhecimento Específico.
Essa distribuição tem moda no intervalo ( 50 ; 60 ] com um quinto dos participantes, seguido dos intervalos ( 70 ; 80 ] e ( 60 ; 70 ].
Caracteriza -se como um máximo local, com cerca de 11 \% do total, a classe de estudantes que deixaram a questão em branco. 6.3.2.4 Comentários sobre as respostas à Questão Discursiva 4 A questão discursiva 4 apresentava a figura de um mapa visual seguindo o modelo processual de projeto, Double Diamond, desenvolvido pelo British Design Council em 2005, que sugere que o processo de design deve compreender quatro fases distintas : descobrir, definir, desenvolver e distribuir.
Abaixo da figura, um texto extraído da página do British Design Council explicava essa figura, definindo o modelo e o conceito de processo criativo subjacente.
O comando pedia que o estudante descrevesse o que ocorre em cada uma das etapas do Double Diamond para chegar a uma boa solução na resolução de um problema de design.
O enunciado da questão estava apresentado de forma clara, com linguagem objetiva e apropriada.
O conteúdo abordado foi adequado em relação às Diretrizes, pois se refere à metodologia do design, conhecimento fundamental para a formação profissional do designer.
O nível de conhecimento exigido é compatível com os cursos de graduação da área.
Embora se refira a um modelo de metodologia específico, a formulação possibilitou a ocorrência de interpretações variadas.
A definição do modelo processual, seguido da ilustração, contribuiu para a clareza e a objetividade do enunciado, mas ao mesmo tempo, levou a respostas que, na descrição das etapas, reproduziram partes do enunciado ou descreveram o que estava apresentado no gráfico, sem demonstrar conhecimento aprofundado do modelo.
Percebeu -se que alguns respondentes conhecem superficialmente o modelo Double Diamond e, a partir do gráfico e da descrição exposta, procuraram fazer deduções que estavam provavelmente baseadas no conhecimento de outros modelos de metodologia existentes, para elaborar as respostas.
O modelo processual apresentado se aplica a qualquer projeto de design, e o padrão de resposta garantiu que estudantes de áreas diferentes do design pudessem responder à questão.
A maior parte das respostas fez uma abordagem genérica com descrições superficiais das etapas, surgindo termos de etapas semelhantes dos demais modelos processuais de projeto em design.
O nível de especificidade que caracteriza o modelo exposto, conforme explicitado no padrão de resposta, foi observado num percentual menor de respostas.
Para a fase ‘ Descobrir ’ foram utilizados pelos respondentes os termos"pesquisa contextualizada"ou"pesquisa/investigação no contexto/situação do problema", descrevendo de forma análoga o termo"imersão".
Para a fase ‘ Definir ’, muitos respondentes mencionaram"definição de um foco de projeto", e alguns mencionaram"definição de um briefing de projeto".
Já na fase ‘ Desenvolver ’, houve maior incidência do termo"geração de alternativas", e poucas respostas comentaram o caráter de"testagem e/ou iteração".
A fase ‘ Distribuir ’ apresentou uma baixa incidência de respostas adequadas.
Muitos respondentes interpretaram essa fase como uma"distribuição aos usuários para testagem", ação que deveria ser mencionada na etapa anterior.
O resultado final mostra um índice elevado de respostas medianas e boas com um índice baixo de respostas fracas.
Dessa forma, pode -se observar que o desempenho dos respondentes foi satisfatório.
As respostas fracas ( notas de 0 a 30 ) foram aquelas em que o respondente apenas citou a fase, mas não fez descrição ou fez a descrição superficialmente, sem citar o nome da fase ( 5,1 \% ) ; as respostas medianas ( notas de 35 a 70 ), em sua maioria, foram aquelas em que o respondente citou as quatro etapas e descreveu adequadamente as três primeiras ( 52,1 \% ) ; e as boas respostas ( notas de 75 a 100 ) foram aquelas em que o respondente citou as quatro fases, descrevendo -as adequadamente ou citou as quatro fases e descreveu adequadamente as três primeiras e superficialmente a última ( 32 \% ).
Deixaram a questão em branco 10,8 \% dos estudantes.
A percepção geral foi a de que os respondentes utilizaram termos de diversos modelos processuais de projeto para descrever as fases do modelo Double Diamond, tentando adequar para as fases como estavam esquematizadas na imagem que acompanhava o enunciado.
A descrição da fase ‘ Distribuir ’ ( Delivery ) pode ter sido prejudicada por uma tradução pouco usual para essa etapa.
Em outras traduções, nota -se o emprego do termo “ Entregar ”,"Entrega"ou"implementação"como tradução dessa fase.
Pode -se supor que o baixo índice de descrição correta dessa fase pode estar relacionado ao fato de que, enquanto aluno, o respondente não vivenciou no seu curso a finalização, a produção e o lançamento do objeto de projeto.
Chamou a atenção, ainda, dentre todas as respostas avaliadas, um respondente ter desenhado o gráfico apresentado no enunciado como base para responder à questão. 6.3.2.5 Análise da Questão Discursiva 5 do Componente de Conhecimento Específico Na Tabela 6.20, constam as informações relativas à questão 5 do conjunto do Componente de Conhecimento Específico.
O desempenho dos estudantes de todo o Brasil nesta questão foi baixo, um pouco superior ao da questão 3.
A nota Média dos estudantes de todo o Brasil foi 18,3.
A maior Média foi registrada na região Centro-Oeste ( 32,4 ), enquanto a menor Média foi registrada na região Norte ( 13,7 ).
Quanto à variabilidade das notas, o Desvio padrão dos alunos do Brasil, como um todo, foi 28,1.
Enquanto o maior desvio foi encontrado na região Centro-Oeste ( 36,4 ), o menor foi encontrado na região Norte ( 22,0 ).
Com exceção da região Centro-Oeste, cuja Mediana foi 30,0, a Mediana, para o Brasil e nas demais regiões, foi zero, associada ao excessivo número de notas zero, como mostrado no Gráfico 6.13, indicando que mais de metade dos alunos receberam essa nota, em cada uma dessas regiões.
Para o conjunto de alunos de Design do Brasil, a nota Máxima foi 100,0, obtida em todas as regiões, com exceção da região Norte, que teve nota Máxima 90,0.
A nota Mínima do Brasil foi zero, o mesmo valor encontrado em todas as regiões. enunciado.
Esse recurso demonstrou que, apesar de não utilizarem exclusivamente os termos relacionados ao modelo em questão souberam se expressar e refletir sobre as fases processuais.
A descrição da fase ‘ Distribuir ’ ( Delivery ) pode ter sido prejudicada por uma tradução pouco usual para essa etapa.
Em outras traduções, nota -se o emprego do termo “ Entregar ”,"Entrega"ou"implementação"como tradução dessa fase.
Pode -se supor que o baixo índice de descrição correta dessa fase pode estar relacionado ao fato de que, enquanto aluno, o respondente não vivenciou no seu curso a finalização, a produção e o lançamento do objeto de projeto.
Chamou a atenção, ainda, dentre todas as respostas avaliadas, um respondente ter desenhado o gráfico apresentado no enunciado como base para responder à questão. 6.3.2.5 Análise da Questão Discursiva 5 do Componente de Conhecimento Específico Na Tabela 6.20, constam as informações relativas à questão 5 do conjunto do Componente de Conhecimento Específico.
O desempenho dos estudantes de todo o Brasil nesta questão foi baixo, um pouco superior ao da questão 3.
A nota Média dos estudantes de todo o Brasil foi 18,3.
A maior Média foi registrada na região Centro-Oeste ( 32,4 ), enquanto a menor Média foi registrada na região Norte ( 13,7 ).
Quanto à variabilidade das notas, o Desvio padrão dos alunos do Brasil, como um todo, foi 28,1.
Enquanto o maior desvio foi encontrado na região Centro-Oeste ( 36,4 ), o menor foi encontrado na região Norte ( 22,0 ).
Com exceção da região Centro-Oeste, cuja Mediana foi 30,0, a Mediana, para o Brasil e nas demais regiões, foi zero, associada ao excessivo número de notas zero, como mostrado no Gráfico 6.13, indicando que mais de metade dos alunos receberam essa nota, em cada uma dessas regiões.
Para o conjunto de alunos de Design do Brasil, a nota Máxima foi 100,0, obtida em todas as regiões, com exceção da região Norte, que teve nota Máxima 90,0.
A nota Mínima do Brasil foi zero, o mesmo valor encontrado em todas as regiões.
No Gráfico 6.13, é representada a distribuição das notas da questão discursiva 5 do Componente de Conhecimento Específico.
Essa distribuição, como ocorreu com a discursiva 3, também tem como frequência modal o intervalo correspondente aos que tiraram nota zero ( 43,8 \% ) que, somados àqueles que deixaram a questão em branco, chegam próximos de dois terços do total.
O próximo intervalo de maior frequência ( 20 ; 30 ] representa um máximo local. 6.3.2.6 Comentários sobre as respostas à Questão Discursiva 5 A questão discursiva 5 iniciava com um texto extraído do livro de Werthein “ A sociedade da informação e seus desafios ” em que o autor alerta para o risco de que o foco sobre a tecnologia possa alimentar a visão ingênua de determinismo tecnológico que percebe a tecnologia como neutra e sem a interferência de fatores sociais e políticos.
O comando solicitava do estudante que, a partir do texto, discorresse sobre a influência da interação entre transformações tecnológicas e fatores socioculturais na atuação profissional do designer.
O enunciado estava claro, com linguagem apropriada.
O conteúdo abordado foi adequado em relação às Diretrizes, apresentando uma questão com nível de complexidade alto, que exigiu uma leitura mais atenta e argumentação mais elaborada.
O enunciado possibilitou a ocorrência de interpretações variadas.
O comando da questão deixou em aberto a abordagem da resposta, não apresentando alguns aspectos concernentes ao assunto, que poderiam contribuir para maior clareza quanto às possibilidades de interpretação.
Tendo em vista que o comando da questão deixou em aberto a abordagem da resposta, a diversidade de interpretações foi significativa.
Na avaliação, foi preciso fazer uma análise atenta para cotejar as respostas com os direcionamentos expostos no padrão de resposta.
Considerando -se que a questão demandou do respondente uma capacidade argumentativa e posicionamento crítico, percebeu -se que a redação precária dificultou a explanação do tópico abordado, prejudicando a resposta.
Percebeu -se que o padrão de resposta se ateve à interpretação da citação apresentada no enunciado e não ao comando de"discorrer sobre a influência das transformações tecnológicas e fatores socioculturais na atuação profissional do designer".
Pode -se destacar uma incidência de dois tipos de respostas : as que se referiram ao perfil do designer como projetista, e as que se referiram à atuação do designer ainda na sua formação acadêmica.
No entanto, as argumentações não alcançaram o que era esperado no padrão de resposta.
A dificuldade na redação e na argumentação prejudicou a dissertação sobre o aspecto destacado pelo respondente para relacionar a influência das transformações tecnológicas e fatores socioculturais na atuação profissional do designer.
O resultado final demonstra um alto índice de respostas consideradas fracas e um baixo índice de respostas consideradas boas, fazendo constatar que o desempenho dos respondentes foi insatisfatório.
Percebeu -se que a grande maioria das respostas ( 57,7 \% ) foi considerada fraca ( notas de 0 a 30 ).
Nelas os respondentes apenas citaram um aspecto da questão e não desenvolveram de forma consistente a argumentação, ficando muitas vezes com uma argumentação incompleta.
As respostas consideradas medianas ( 15,7 \% ) e as respostas consideradas boas ( 6,6 \% ) apresentaram, pelo menos, um aspecto da questão, mas possuíam um grau de argumentação mais elevado.
A diferenciação entre elas se deu pela qualidade da argumentação e pela existência de uma conclusão da argumentação nas respostas consideradas boas.
As Respostas Medianas correspondem às notas de 35 a 70 e as Respostas Boas, às notas de 75 a 100.
Deixaram a Questão 5 em branco um total de 20 \% dos respondentes.
Esse resultado permite supor que os respondentes não receberam em sua formação elementos que possibilitassem a reflexão crítica.
Além disso, percebeu -se a dificuldade de se expressar por escrito, o que dificultou a argumentação e a articulação dos conceitos que o respondente apresentou no seu texto.
Esse fato exigiu do corretor uma leitura atenta e um tempo de correção mais lento, com o objetivo de identificar, na redação, conteúdos que se relacionassem com os critérios do padrão de resposta. 6.3.3 Considerações Finais O resultado da avaliação do Enade/2018, sob a perspectiva das questões discursivas específicas, reveste -se de uma importância particular por fornecer dados e informações para uma análise de como está a qualidade da formação dos estudantes brasileiros de design em relação a conteúdos concernentes a : História do Design Brasileiro ; métodos processuais para o desenvolvimento de Projeto em Design ; e posicionamento crítico relativo à atuação do designer enquanto projetista.
A seleção das temáticas possibilita um olhar diverso para a formação do designer e a \textbf{observação} de aspectos que se complementam na atuação profissional.
Percebeu -se que o recurso de ilustração nas questões facilita a compreensão do enunciado.
No entanto, nota -se que muitos respondentes da questão relativa ao estilo gráfico do designer Alexandre Wollner se detiveram apenas na descrição da imagem e não na descrição do estilo, o que denota deficiência de conhecimento.
Essa questão permite avaliar tanto o nível de conhecimento da História do Design Brasileiro quanto o modo como os respondentes utilizam a nomenclatura para descrever o estilo gráfico.
A imagem do modelo processual Double Diamond também contribuiu para o entendimento do enunciado, possibilitando que os respondentes o associassem a outros modelos processuais e apresentassem definições compatíveis.
O resultado demonstra que há um conhecimento formado com relativa consistência sobre conceitos pertinentes à metodologia de projeto.
Por outro lado, evidencia -se uma deficiência acentuada na formação relativa à reflexão sobre a atuação do designer, enquanto projetista, bem como sobre sua relação com o desenvolvimento tecnológico e as questões sociais.

% matched lemmas: aprendizado, centrado, destruir, esporádico, mentira, observação, personificar, peso, sócio-histórico-cultural, étnico-racial
\end{textsample}
